\section{结论}
\label{sec:conclusion}

本文提出并在一个经过完整质量审计的批次上验证了一个面向基因编辑影响因素发现的科学发现平台，以 DeepCRISPR 公开数据作为案例研究。平台的定位不是构造单一最优预测器，也不是证明表观遗传环境必然提升预测，而是把质量控制、数据划分与泄漏审计、模型归因、统计推断、环境增量分析、序列模式发现、细胞背景比较与证据整合组织为统一且可追溯的流程。

案例研究的主要结论如下，每条均标注其证据强度：

\begin{enumerate}
\item \textbf{数据可信度（穷举核对，强）}：1\,344 次运行的实验矩阵完整，1\,344/1\,344 的 train$\cap$test 序列重叠与反向互补重叠为 0，划分摘要逐条复核一致，全批运行于同一数据、代码与数值栈之上（表~\ref{tab:dataquality}）。
\item \textbf{泄漏整改改变结论方向（对照观测，中）}：在排除跨细胞系的同源序列后，\emph{mixed} 中位 $R^{2}$ 由 0.118 降至 0.073，LOCO 由 $+0.036$ 降至 $-0.008$，而 \emph{single} 划分基本不变。这一差异化变化表明此前较高的跨细胞系性能至少有相当部分由同源序列重复支撑。
\item \textbf{预测能力有限且跨细胞系泛化失败（定量观测，强）}：细胞系内中位 $R^{2}$ 为 0.067--0.120，混合划分为 0.041--0.094，而真实 LOCO 的 7 个模型配置全部落在 $-0.112$ 至 $+0.010$，并呈显著细胞系异质性（HEK293T 最低，$R^{2}=-0.266$）。因此本文不主张一个具有跨细胞系泛化能力的通用预测器，而将平台定位为\textbf{受控、可追溯的影响因素研究框架}。
\item \textbf{序列信号集中在 PAM 邻近位置（多模型描述性汇总，中）}：跨模型平均归因谱的峰值位置为第 18 位（其后为 20、23、19），XGBoost 与 MLP 分别有 91.7\% 与 86.5\% 的运行把最高归因位置落在 PAM 邻近窗口（17--20）。但按粗粒度区域汇总时，只有 XGBoost 与 Transformer 的最高区域为 PAM 邻近种子区，CNN 与 MLP 更偏向种子核心区——本文据此限定表述，不使用"所有模型类一致指向 PAM 邻近区"的说法。
\item \textbf{环境因素未获得增量预测证据（区间与门槛判定，强）}：CTCF、DNase、H3K4me3、RRBS 的效应量为 $3\times10^{-4}$--$1.5\times10^{-3}$，比预设门槛 0.01 低 1--2 个数量级；跨模型区间全部跨 0；区组析因方差分析四个主效应 $p=0.975$--$0.9996$。四者在证据框架中均为 \textbf{Inconclusive}。这一结论是\emph{未检出}，不是\emph{无作用}。
\item \textbf{感受野带来小幅但可复现的提升（配对实验，中）}：$k=7$ 相对 $k=3$ 提高 $+0.0169$（区间 $+0.0077$--$+0.0257$，79.7\% 配对为正），而 $k=5$ 相对 $k=3$ 的差异不可区分于 0（$+0.0046$，区间跨 0）。候选模式长度不随核大小改变，说明卷积核是感受野参数而非模式长度。
\item \textbf{可检验的候选假设（模型内部反事实，弱—中）}：平台输出第 18 位碱基替换与 PAM 邻近候选模式破坏两类最小扰动候选，用于后续实验验证。这些候选基于模型归因与观测性碱基关联，\textbf{不构成因果证据}。
\end{enumerate}

综上，本文的贡献不在于报告一个高精度预测器，而在于展示：在严格控制数据真实性、明确划分身份并审计泄漏之后，仍可以从有限预测能力中提取可解释、可检验的科学信号，同时明确哪些问题当前数据无法回答。后续工作有两条主线：一是在独立编辑数据集上检验流程与候选的可迁移性；二是对本文提出的候选假设开展最小扰动实验验证。
